% from ega4-14-fr.tex
\oldpage[IV-3]{220}

\emph{I) Réduction au cas où $Y$ est intègre.} -- Si l'on a prouvé la proposition
pour chacun des sous-préschémas réduits $Y_i$ ayant pour espace sous-jacent une
composante irréductible de $Y$, et pour les images réciproques $f^{-1}(Y_i)$, il est
clair que le préschéma $Y'$ \emph{somme} des $Y'_i$ correspondants répondra à la
question. On peut donc supposer $Y$ \emph{intègre} et on va dans ce qui suit se
borner à ce cas. Alors, dans la conclusion, on pourra aussi prendre $Y'$ intègre
(quitte à le remplacer par une composante irréductible convenable).

\emph{II) Caractère local sur $Y$.} -- Nous allons montrer que si l'on peut
recouvrir $Y$ par des ouverts $U_j$ en nombre fini, tels que, pour tout $j$, la
conclusion de la proposition soit vraie pour le morphisme $f^{-1}(U_j)\to U_j$,
restriction de $f$, alors la conclusion est vraie également pour $f$. En effet, on
peut évidemment supposer les $U_j$ affines, de sorte que $U_j=\operatorname{Spec}(A_j)$, où
$A_j$ est un anneau intègre noethérien dont le corps des fractions $K=R(Y)$ est le
corps des fonctions rationnelles sur $Y$. Pour tout $j$, il y a par hypothèse une
$A_j$-algèbre finie intègre $A'_j$ telle que l'homomorphisme $A_j\to A'_j$ soit
injectif (\textbf{I}, 1.2.7) et que le morphisme correspondant
$g_j:U'_j=\operatorname{Spec}(A'_j)\to\operatorname{Spec}(A_j)=U_j$ vérifie les conditions de la proposition
(pour $U_j$ et $f^{-1}(U_j)$). Soit alors $K'$ une extension finie de $K$ contenant
les corps des fractions de tous les $A'_j$ (qui sont des extensions finies de $K$).
Considérons le normalisé $Y''$ de $Y$ dans $K'$ (\textbf{II}, 6.3.8), qui est de la
forme $\operatorname{Spec}(\mathcal B)$, où $\mathcal B$ est une $\mathcal O_Y$-Algèbre entière
quasi-cohérente, fermeture intégrale de $\mathcal O_Y$ dans $K'$ (\textbf{II},
6.3.4). Ces définitions prouvent que pour tout $j$,
\[
  A'_j=\Gamma(U'_j,\mathcal O_{U'_j})
      =\Gamma(U_j,(g_j)_*(\mathcal O_{U'_j}))
\]
s'identifie à une sous-$A_j$-algèbre finie de $\Gamma(U_j,\mathcal B)$~; autrement
dit, $(g_j)_*(\mathcal O_{U'_j})=\mathcal A'_j$ est une
$\mathcal O_{U_j}$-Algèbre cohérente, sous-Algèbre de $\mathcal B|U_j$. Il existe
donc un sous-$\mathcal O_Y$-Module $\mathcal C'_j$ cohérent de $\mathcal B$ tel que
$\mathcal C'_j|U_j=\mathcal A'_j$ (\textbf{I}, 9.4.7). Si l'on pose
$\mathcal C'=\sum_j\mathcal C'_j$, la sous-$\mathcal O_Y$-Algèbre $\mathcal B'$ de
$\mathcal B$ engendrée par $\mathcal C'$ est cohérente puisque $\mathcal B$ est une
$\mathcal O_Y$-Algèbre entière. Montrons alors que $Y'=\operatorname{Spec}(\mathcal B')$ répond à
la question. En effet, il est clair que le morphisme $g:Y'\to Y$ est fini surjectif
et que $Y'$ est intègre~; en outre, pour tout $j$, $\Gamma(U_j,\mathcal B')$ est une
$A'_j$-algèbre finie, autrement dit le morphisme $g^{-1}(U_j)\to U_j$, restriction
de $g$, se factorise en $g^{-1}(U_j)\to U'_j\to U_j$, et comme l'existence locale
de sections est stable par changement de base, cela établit notre assertion, tout
$y\in Y$ appartenant à un $U_j$.

\emph{III) Réduction au cas où $Y$ est intègre, local et géométriquement
unibranche.} -- Supposons d'abord que la proposition ait été prouvée lorsque $Y$
est intègre et \emph{local} (avec $Y'$ intègre), et montrons qu'elle est valable
lorsque $Y$ est intègre (noethérien) affine quelconque. En effet, en vertu de la
réduction II), il suffit de prouver que pour tout point $y\in Y$, la proposition est
vraie pour un voisinage ouvert affine $V$ de $y$ dans $Y$. Soient $Y=\operatorname{Spec}(A)$,
$Y_1=\operatorname{Spec}(A_{\mathfrak p})$, où $\mathfrak p=\mathfrak j_y$, et posons
$X_1=X\times_Y Y_1$~; par hypothèse, il existe un morphisme fini surjectif
$g_1:Y'_1\to Y_1$, où $Y'_1=\operatorname{Spec}(B_1)$, $B_1$ étant une
$A_{\mathfrak p}$-algèbre intègre finie, donc un anneau semi-local, tel que $g_1$
vérifie les conditions de l'énoncé pour $Y_1$ et $X_1$. Si $y'_j$
($1\leq j\leq r$) sont les points fermés de $Y'_1$, il y a donc un recouvrement de
$Y'_1$ par des ouverts $U'_j$ tels que $y'_j\in U'_j$ et qu'il existe une
$U'_j$-section $h'_j$ de $X\times_Y U'_j$ ($1\leq j\leq r$). Le
$A_{\mathfrak p}$-module $B_1$ admet un système fini de générateurs de la forme
$z_i/s$ (avec $s\in A-\mathfrak p$,

